
\def\PUDcodigo{ACE003}

\if\COMPILAR0
    \def\PUDcodacad{04507.2}
\else
    \def\PUDcodacad{04506.2}
\fi

\def\PUDdisciplina{Cálculo 1}

\def\PUDsemestre{S1}

\def\PUDhoras{80}

%NOVOS ---------------------------------------------------
\def\PUDhorasTeoria{80}
\def\PUDhorasPratica{0}

\def\PUDinterECA{---}

\def\PUDatuaisECA{---}

\def\PUDrecursos{Projetor multimídia; Quadro branco e pincel.}

%----------------------------------------------------------

\def\PUDcreditos{4}

\def\PUDprerequisito{---}

\def\PUDpreacad{---}

\def\PUDobjetivos{Conhecer as ferramentas básicas do Cálculo Diferencial e Integral I, bem como se capacitar a aplicar tais ferramentas na resolução de problemas afins à atividade profissional.}

\def\PUDementa{Noções preliminares de cálculo;  Limites e continuidade de funções;  Derivação;  Aplicações da derivada;  Integração;  Aplicações da integral.}

\def\PUDconteudo{ &  Noções preliminares: Números reais;  Plano cartesiano;  Conceito de função;  Tipologia das funções;  Composição e inversão de funções; 
\\ 
 &  Limites e continuidade de funções: Noção intuitiva e exemplos;  Definição de limite;  Propriedades operatórias dos limites;  Teoremas sobre limites;  Limites laterais;  Limites fundamentais;  Funções contínuas; 
\\ 
 &  Derivação: Velocidade;  Coeficiente angular;  Definição de derivada;  Função derivada;  Propriedades operatórias da derivada;  Derivadas das funções elementares;  Regra da cadeia;  Derivada da função inversa;  Derivação implícita; 
\\ 
 &  Aplicações da derivada: Estudo da variação das funções;  Funções convexas;  Máximos e mínimos;  Taxas de variação;  Taxas de variação relacionadas;  Expressões indeterminadas (regra de L’Hopital); 
\\ 
 &  Introdução a integração. }

\def\PUDmetodo{Aulas expositórias, em que busca-se sempre mostrar aplicações da disciplina, de forma a servir como estopim para a curiosidade do aluno, motivando-o a estudar e aprofundar-se no mundo dessa subárea da Matemática. Além disso, busca-se, sempre que possível, envolver a monitoria na metodologia de ensino, de tal forma que melhore o aprendizado do aluno, amenizando assim suas dificuldades e os ajudando a progredir no curso. 
%Aulas expositórias que podem ser teóricas e/ou práticas, onde as práticas no laboratório serão marcadas no decorrer da disciplina.
}

\def\PUDavalia{A avaliação será feita com aplicação de provas teóricas, além da possibilidade de inclusão de trabalhos, seminários e listas de exercícios no decorrer da disciplina.
%A avaliação será feita com aplicação de provas teóricas e/ou práticas, além da possibilidade de inclusão de trabalhos e seminários no decorrer da disciplina.
}

\def\PUDbibbasica{ & \href{www.biblioteca.ifce.edu.br/index.asp?codigo_sophia=23696}{Flemming}, Diva Marília.  Cálculo A: funções, limite, derivação, integração.  6º ed.  São Paulo, SP : Pearson Prentice Hall, 2007.  449p.  ISBN: 857605115-X. 
\\ 
 &  \href{www.biblioteca.ifce.edu.br/index.asp?codigo_sophia=63211}{Stewart}, James.  Cálculo 1.  São Paulo, SP : Cengage Learning, 2016.  525p.  ISBN: 9788522112586.
\\ 
 &  \href{www.biblioteca.ifce.edu.br/index.asp?codigo_sophia=23890}{Thomas}, George B.  Cálculo 1.  11º ed.  São Paulo, SP: Addison-Wesley, 2009.  ISBN: 9788588639317. }

\def\PUDbibcomplementar{ &  \href{www.biblioteca.ifce.edu.br/index.asp?codigo_sophia=23737}{Guidorizzi}, Hamilton Luiz.  Um curso de cálculo - v. 1.  5º ed.  Rio de Janeiro, RJ : LTC, 2001.  ISBN: 9788521612599. 
\\ 
 &  \href{www.biblioteca.ifce.edu.br/index.asp?codigo_sophia=24328}{Leithold}, Louis.  O cálculo com geometria analítica.  v.  1.  3º ed.  São Paulo, SP : Harbra, 1994.  ISBN: 8529400941.  
\\ 
 &  \href{www.biblioteca.ifce.edu.br/index.asp?codigo_sophia=61965}{Hoffmann}, Laurence.  Cálculo: um curso moderno e suas aplicações.  11º ed.  Rio de Janeiro, RJ : LTC, 2016.  661p.  ISBN: 9788521625315.
\\
 &  \href{www.biblioteca.ifce.edu.br/index.asp?codigo_sophia=40002}{Anton}, Howard.  Cálculo 1.  8º ed.  Porto Alegre, RS. Bookman, 2009.  680p.  ISBN: 9788560031634.
\\
 &  \href{www.biblioteca.ifce.edu.br/index.asp?codigo_sophia=62611}{Hughes-Hallett}, Deborah.  Cálculo aplicado.  4º ed.  Rio de Janeiro, RJ: LTC, 2012.  483p.  ISBN: 9788521620518. }

\def\PUDautor{Samuel Vieira Dias}

\def\PUDdata{2013-04-23}

\def\PUDrevisao{1}

\def\PUDrevisor{Samuel Vieira Dias}

\def\PUDdatarevisao{2017-05-09}

\PUDcorpo
